\begin{historical}
The concept of vacuum has been contested since antiquity. Aristotle rejected the possibility of void—nature abhors a vacuum—and medieval natural philosophy largely followed. Torricelli's 1643 experiment demonstrated that a sustained void could exist above a column of mercury, but the space was considered merely empty of air, not of all physical influence. Newton's mechanics required absolute space as a reference frame but assigned no physical properties to empty regions. The nineteenth-century luminiferous aether filled the vacuum with an elastic medium to carry light, but the Michelson-Morley experiment (1887) and Einstein's special relativity (1905) eliminated the aether without replacing it. Quantum mechanics reopened the question. Dirac's 1930 electron theory predicted that the vacuum teemed with an infinite sea of negative-energy electrons—the Dirac sea—whose holes appeared as positrons. Casimir's 1948 prediction that two uncharged conducting plates would attract due to restricted vacuum modes, confirmed experimentally by Lamoreaux in 1997, established that the quantum vacuum exerts measurable forces.

By the time quantum field theory matured on curved backgrounds, the vacuum had already shed any pretence of emptiness. In 1966–69, Leonard Parker showed that the expansion of the universe itself could create particles from the vacuum, the first rigorous demonstration that a time-dependent spacetime geometry converts vacuum fluctuations into real excitations. In 1973, Stephen Fulling proved that the definition of \QSOpen{}particle\QSClose{} depends on the observer's choice of quantisation, while Paul C. W. Davies and William G. Unruh showed that a uniformly accelerated observer in flat spacetime perceives a thermal bath where an inertial observer sees none. In 1974, Stephen Hawking extended these ideas to black holes, deriving that they emit thermal radiation at a temperature inversely proportional to their mass. The result stunned even Hawking—he had set out to disprove it. Jacob Bekenstein had already argued on information-theoretic grounds that black holes must carry entropy proportional to their horizon area, but without a temperature the claim seemed formal. Hawking's calculation supplied the temperature, and with it the entropy became physical: $S = k_B c^3 A / (4G\hbar)$. The four laws of black hole mechanics, catalogued by Bardeen, Carter, and Hawking in 1973 as a mathematical analogy to thermodynamics, turned out to be thermodynamics.

The implications extended beyond black holes. If particle content depends on the observer, then so does the vacuum energy density that drives cosmological expansion, complicating the already severe cosmological constant problem. The Schwinger effect (1951), in which a sufficiently strong electric field pulls electron-positron pairs from the vacuum, provided an independent line of evidence that the vacuum is a reservoir from which real particles can be extracted given enough energy. Birrell and Davies synthesised these threads in their 1982 monograph \textit{Quantum Fields in Curved Space}, which remains the standard reference. By the late twentieth century, the vacuum had been transformed from the simplest concept in physics—nothing—into one of the most contested.
\end{historical}
